\bta{2005}


\begin{enumerate}
	%\renewcommand{\labelenumi}{\arabic{enumi}.}
	% A(\Alph) a(\alph) I(\Roman) i(\roman) 1(\arabic)
	%设定全局标号series=example	%引用全局变量resume=example
	%[topsep=-0.3em,parsep=-0.3em,itemsep=-0.3em,partopsep=-0.3em]
	%可使用leftmargin调整列表环境左边的空白长度 [leftmargin=0em]
	\item
	填空题 (本题共 5 小题，每小题 5 分，满分 25 分)
	\begin{enumerate}
		%\renewcommand{\labelenumi}{\arabic{enumi}.}
		% A(\Alph) a(\alph) I(\Roman) i(\roman) 1(\arabic)
		%设定全局标号series=example	%引用全局变量resume=example
		%[topsep=-0.3em,parsep=-0.3em,itemsep=-0.3em,partopsep=-0.3em]
		%可使用leftmargin调整列表环境左边的空白长度 [leftmargin=0em]
		\item
		已知 $f^{\prime}\left(x_{0}\right)=3$, 则 $\lim\limits _{x \rightarrow 0} \frac{f\left(x_{0}\right)-f\left(x_{0}-2 x\right)}{x}=$ \tk{$ 6 $} 。
		
		
		\item
		 设 $f(x)$ 的一个原函数为 $e^{x^{2}}$, 则 $\int x f^{\prime}(x) d x=$ \tk{$ (2x^{2}-1)e^{x^{2}} +C $} 。
		
		
		
		
		\item
		数量场 $u=x^{2}+2 y^{2}+3 z^{2}$ 在点 $(1,1,-1)$ 的最大方向微商值为 \tk{$ 2\sqrt{14} $} 。
		
		
		\item
		级数 $\sum\limits_{n=1}^{\infty} \frac{x^{2 n}}{2^{n}+3^{n}}$ 的收敛半径为 \tk{$\sqrt{3}$} 。
		
		
		\item
		微分方程 $y^{\prime}-\frac{y}{x}=1$ 的通解为 \tk{$ y=x(C+ \ln |x|) $} 。
		
		
	\end{enumerate}

\banswer{
\begin{enumerate}
	%\renewcommand{\labelenumi}{\arabic{enumi}.}
	\item
	$6$。凑成导数定义的形式：
	\begin{align*}
	\lim_{x\to0}\frac{f\left(x_{0}\right)-f\left(x_{0}-2x\right)}{x}
	&=2\lim_{x\to0}\frac{f\left(x_{0}-2x\right)-f\left(x_{0}\right)}{-2x}\\
	&=2f^{\prime}\left(x_{0}\right)=6 .
	\end{align*}

	\item
	$\left(2x^{2}-1\right)e^{x^{2}}+C$。由 $e^{x^{2}}$ 是 $f$ 的一个原函数得
	$f(x)=\left(e^{x^{2}}\right)^{\prime}=2xe^{x^{2}}$，且 $\int f(x)dx=e^{x^{2}}+C$。分部积分：
	\begin{align*}
	\int xf^{\prime}(x)dx&=xf(x)-\int f(x)dx\\
	&=2x^{2}e^{x^{2}}-e^{x^{2}}+C=\left(2x^{2}-1\right)e^{x^{2}}+C .
	\end{align*}

	\item
	$2\sqrt{14}$。最大方向微商（方向导数）的值等于梯度的模。
	$\operatorname{grad}u=(2x,4y,6z)$，在 $(1,1,-1)$ 处为 $(2,4,-6)$，故最大值为
	\[
	\left|\operatorname{grad}u\right|=\sqrt{2^{2}+4^{2}+(-6)^{2}}=\sqrt{56}=2\sqrt{14} ,
	\]
	方向为 $(1,2,-3)$ 的单位向量方向。

	\item
	$\sqrt{3}$。令 $t=x^{2}$，级数化为 $\sum\limits_{n=1}^{\infty}\frac{t^{n}}{2^{n}+3^{n}}$。由
	\[
	\sqrt[n]{\frac{1}{2^{n}+3^{n}}}=\frac{1}{3}\sqrt[n]{\frac{1}{1+(2/3)^{n}}}\longrightarrow\frac13
	\]
	知它关于 $t$ 的收敛半径为 $3$，即 $x^{2}<3$。故原级数当 $|x|<\sqrt3$ 时收敛、
	$|x|>\sqrt3$ 时发散，收敛半径 $R=\sqrt3$。

	\item
	$y=x(C+\ln|x|)$。这是一阶线性方程 $y^{\prime}-\frac1xy=1$，取积分因子 $\frac1x$，得
	\[
	\left(\frac{y}{x}\right)^{\prime}=\frac{xy^{\prime}-y}{x^{2}}=\frac1x ,
	\]
	积分得 $\frac{y}{x}=\ln|x|+C$，即 $y=x(C+\ln|x|)$。

	\par\vspace{1mm}\noindent
	\textbf{注：}{\kaishu
	原题面印作 $y^{\prime}+\frac1x=1$，其中根本不含未知函数 $y$，
	直接积分只得 $y=x-\ln|x|+C$，与所给答案 $y=x(C+\ln|x|)$ 不符。
	由答案反推，原方程应为 $y^{\prime}-\frac{y}{x}=1$（已代入验证满足），故据此更正题面。
	}\par
\end{enumerate}
}

\item
单项选择题 (本题共 5 小题，每小题 5 分，满分 25 分)
\begin{enumerate}
	%\renewcommand{\labelenumi}{\arabic{enumi}.}
	% A(\Alph) a(\alph) I(\Roman) i(\roman) 1(\arabic)
	%设定全局标号series=example	%引用全局变量resume=example
	%[topsep=-0.3em,parsep=-0.3em,itemsep=-0.3em,partopsep=-0.3em]
	%可使用leftmargin调整列表环境左边的空白长度 [leftmargin=0em]
	\item
	设 $f(0)=0$, 则 $f(x)$ 在点 $x=0$ 可导的充要条件为 \xzanswer{C} 
	
	
\fourchoices
{$\lim\limits _{t \rightarrow 0} \frac{1}{t^{2}} f\left(t^{2}\right)$ 存在}
{$\lim\limits _{t \rightarrow 0} \frac{1}{t^{2}} f(t-\sin t)$ 存在}
{$\lim\limits _{t \rightarrow 0} \frac{1}{t} f(\ln (1+t))$ 存在}
{$\lim\limits _{t \rightarrow 0} \frac{1}{t^{2}}[f(2 t)-f(t)]$ 存在}

\item 
设曲面 $x^{2}+y^{2}-\frac{z^{2}}{4}=1$ 在点 $(1,1,2)$ 处的法线为 $L$, 又设 $L_{1}: \frac{x-1}{2}=\frac{y-2}{1}=$ $\frac{z-1}{0} ; \quad \pi: x+y+4 z=1$, 则 \xzanswer{A} 


\fourchoices
{$L$ 与 $L_{1}$ 相交, 且 $L$ 平行于 $\pi$}
{$L$ 与 $L_{1}$ 相交, 且 $L$ 垂直于 $\pi$}
{$L$ 与 $L_{1}$ 异面, 且 $L$ 平行于 $\pi$}
{$L$ 与 $L_{1}$ 异面, 且 $L$ 垂直于 $\pi$}

\item 
设 $S$ 是柱面 $x^{2}+y^{2}=R^{2}(0 \leqslant z \leqslant R)$ 的外侧, 则 $\iint_{S}\left(x^{2}+y^{2}\right) d x d y$ 的值为 \xzanswer{D} 


\fourchoices
{$2 \pi R^{3}$}
{$2 \pi R^{4}$}
{$\pi R^{4}$}
{$ 0 $}

\item 
设级数 $\sum\limits_{n=1}^{\infty} a_{n}$ 收敛, 则下列结论中正确的是 \xzanswer{B} 


\fourchoices
{级数 $\sum\limits_{n=1}^{\infty} a_{n}$ 收敛}
{级数 $\sum\limits_{n=1}^{\infty} \sqrt[n]{n} a_{n}$ 收敛}
{级数 $\sum\limits_{n=1}^{\infty} \frac{(-1)^{n}}{\sqrt{n}} a_{n}$ 收敛}
{级数 $\sum\limits_{n=1}^{\infty} \frac{a_{n}}{n}$ 收敛}

\item 
设 $f(x)=x-L(0 \leqslant x \leqslant 2 L)$, 则其以 $2 L$ 为周期的傅里叶级数在点 $x=-\frac{L}{2}$ 处收 敛于 \xzanswer{C} 


\fourchoices
{$-\frac{L}{2}$}
{$-\frac{3 L}{2}$}
{$\frac{L}{2}$}
{$\frac{3 L}{2}$}


	
	
\end{enumerate}

\banswer{
\begin{enumerate}
	%\renewcommand{\labelenumi}{\arabic{enumi}.}
	\item
	\textbf{C}。因 $f(0)=0$，$f$ 在 $x=0$ 可导 $\iff$ 双侧极限
	$\lim\limits_{u\to0}\frac{f(u)}{u}$ 存在。
	对 C，令 $u=\ln(1+t)$，当 $t$ 在 $0$ 的两侧变动时 $u$ 与 $t$ 同号且同样取遍 $0$ 的两侧，
	又 $\frac{u}{t}\to1$，于是
	\[
	\frac{f(\ln(1+t))}{t}=\frac{f(u)}{u}\cdot\frac{u}{t} ,
	\]
	其极限存在与 $\lim\limits_{u\to0}\frac{f(u)}{u}$ 存在互为充要，故 C 是充要条件。
	A 中 $t^{2}\to0^{+}$，只刻画右导数；B 中 $t-\sin t\sim\frac{t^{3}}{6}>0$，
	且 $\frac{t-\sin t}{t^{2}}\to0$，取 $f(x)=|x|$ 时 A、B 的极限都存在而 $f$ 不可导，
	故 A、B 必要而不充分。
	D 中若 $f$ 可导且 $f^{\prime}(0)\neq0$（如 $f(x)=x$），
	则 $\frac{f(2t)-f(t)}{t^{2}}\sim\frac{f^{\prime}(0)}{t}\to\infty$，极限不存在，故 D 连必要性也不具备。

	\item
	\textbf{A}。令 $F=x^{2}+y^{2}-\frac{z^{2}}{4}-1$，则
	$\operatorname{grad}F=\left(2x,2y,-\frac z2\right)$，在 $(1,1,2)$ 处为 $(2,2,-1)$，
	即法线 $L$ 过 $(1,1,2)$、方向 $\bm{s}=(2,2,-1)$。
	因 $\bm{s}\cdot(1,1,4)=2+2-4=0$，而 $1+1+4\times2=10\neq1$ 说明 $(1,1,2)\notin\pi$，
	故 $L$ 平行于 $\pi$（不在 $\pi$ 内）。
	又 $L$ 上的点为 $(1+2s,1+2s,2-s)$，$L_{1}$ 上的点为 $(1+2u,2+u,1)$；
	由 $2-s=1$ 得 $s=1$，对应点 $(3,3,1)$，而 $u=1$ 时 $L_{1}$ 上的点也是 $(3,3,1)$，
	故 $L$ 与 $L_{1}$ 相交。选 A。

	\item
	\textbf{D}。柱面 $x^{2}+y^{2}=R^{2}$ 的侧面上每点的法向量都与 $z$ 轴垂直，
	即方向余弦 $\cos\gamma\equiv0$，于是
	\[
	\iint_{S}\left(x^{2}+y^{2}\right)dxdy=\iint_{S}\left(x^{2}+y^{2}\right)\cos\gamma\,dS=0 .
	\]
	（等价地：该曲面在 $xOy$ 面上的投影是一条曲线，面积为零。）

	\item
	\textbf{B}。记 $b_{n}=\sqrt[n]{n}$。当 $n\geqslant3$ 时 $\{b_{n}\}$ 单调下降且
	$1<b_{n}\leqslant\sqrt[3]{3}$，即单调有界。由阿贝尔判别法
	（$\sum a_{n}$ 收敛且 $\{b_{n}\}$ 单调有界 $\Longrightarrow\sum a_{n}b_{n}$ 收敛）
	知 $\sum\sqrt[n]{n}\,a_{n}$ 收敛。
	C 不真：取 $a_{n}=\frac{(-1)^{n}}{\sqrt n}$（由莱布尼茨判别法 $\sum a_{n}$ 收敛），则
	$\sum\frac{(-1)^{n}}{\sqrt n}a_{n}=\sum\frac1n$ 发散。

	\par\vspace{1mm}\noindent
	\textbf{注：}{\kaishu
	本题题面有缺陷。其一，选项 A 与题设条件逐字相同（都是“级数 $\sum a_{n}$ 收敛”），
	显系印刷脱字，按通行版本应为 $\sum\left|a_{n}\right|$（或 $\sum a_{n}^{2}$）收敛。
	其二，选项 D 中 $\left\{\frac1n\right\}$ 同样单调有界，
	由与 B 完全相同的阿贝尔判别法，$\sum\frac{a_{n}}{n}$ 也必收敛
	（数值校验：取 $a_{n}=\frac{(-1)^{n}}{\sqrt n}$ 时部分和稳定于 $-0.76515$），
	故按印出的题面 D 亦为真命题，正确选项不唯一。此处仍按标准答案取 B。
	}\par

	\item
	\textbf{C}。以 $2L$ 为周期延拓后，$x=-\frac L2$ 与 $x=-\frac L2+2L=\frac{3L}{2}$ 是同一点，
	而 $\frac{3L}{2}\in(0,2L)$ 是 $f$ 的连续点（延拓后的间断点只在 $x=2kL$ 处），
	由狄利克雷收敛定理，级数在该点收敛于
	\[
	f\!\left(\frac{3L}{2}\right)=\frac{3L}{2}-L=\frac{L}{2} .
	\]
	也可直接算系数：对 $f(x)=x-L$ 有 $a_{n}=0$、$b_{n}=-\frac{2L}{n\pi}$，故级数为
	$-\frac{2L}{\pi}\sum\limits_{n=1}^{\infty}\frac{1}{n}\sin\frac{n\pi x}{L}$，在 $x=-\frac L2$ 处为
	\[
	\frac{2L}{\pi}\left(1-\frac13+\frac15-\cdots\right)=\frac{2L}{\pi}\cdot\frac{\pi}{4}=\frac{L}{2} .
	\]

	\par\vspace{1mm}\noindent
	\textbf{注：}{\kaishu
	易错之处是把 $-\frac L2$ 直接代入表达式 $x-L$ 得 $-\frac{3L}{2}$（选项 B），
	或误认为 $-\frac L2$ 是间断点而取跳跃平均值。
	实际上 $f$ 只定义在 $[0,2L]$ 上，$-\frac L2$ 处的值由周期延拓给出，
	且该点是连续点，故收敛于 $\frac L2$。
	}\par
\end{enumerate}
}

\newpage
\item
(5 小题，每小题 8 分，满分 40 分)
\begin{enumerate}
	%\renewcommand{\labelenumi}{\arabic{enumi}.}
	% A(\Alph) a(\alph) I(\Roman) i(\roman) 1(\arabic)
	%设定全局标号series=example	%引用全局变量resume=example
	%[topsep=-0.3em,parsep=-0.3em,itemsep=-0.3em,partopsep=-0.3em]
	%可使用leftmargin调整列表环境左边的空白长度 [leftmargin=0em]
	\item
计算极限 $\lim\limits _{x \rightarrow 0} \frac{\sin x-\sin (\sin x)}{x^{3}}$。
	
	
	\item
	计算广义积分 $\int_{0}^{+\infty} \frac{x}{\left(2+x^{2}\right) \sqrt{1+x^{2}}} d x$。
	


\item
 利用欧拉积分计算 $\int_{0}^{1} \frac{x^{6}}{\sqrt[6]{1-x^{6}}} d x$。


\item
$f(u, v)$ 具有二阶连续偏导数, $z=f\left(x y^{2}, \frac{y}{x}\right)$, 求 $\frac{\partial z}{\partial x}, \frac{\partial^{2} z}{\partial x \partial y}$。


\item
计算二重积分 $\int_{0}^{1} d x \int_{\sqrt[3]{x}}^{1} \cos y^{2} d y$。
	
	
	
\end{enumerate}



\banswer{
	\begin{enumerate}
		%\renewcommand{\labelenumi}{\arabic{enumi}.}
		% A(\Alph) a(\alph) I(\Roman) i(\roman) 1(\arabic)
		%设定全局标号series=example	%引用全局变量resume=example
		%[topsep=-0.3em,parsep=-0.3em,itemsep=-0.3em,partopsep=-0.3em]
		%可使用leftmargin调整列表环境左边的空白长度 [leftmargin=0em]
		\item
		$  \frac{ 1 }{ 6 }  $
		\item 
		$ \frac{\pi}{4} $
		\item 
		$ \frac{\pi}{18} $
		\item 
		$\frac{\partial z}{\partial x}=y^{2} \frac{\partial f}{\partial u}-\frac{y}{x^{2}} \frac{\partial f}{\partial v}$，
		$\frac{\partial^{2} z}{\partial x \partial y}=2 y \frac{\partial f}{\partial u}-\frac{1}{x^{2}} \frac{\partial f}{\partial v}+2 x y^{3} \frac{\partial^{2} f}{\partial u^{2}}-\frac{y^{2}}{x} \frac{\partial^{2} f}{\partial u \partial v}-\frac{y}{x^{3}} \frac{\partial^{2} f}{\partial v^{2}}$
		\item 
	$=\frac{1}{2}(\sin 1+\cos 1-1)$
	
	\end{enumerate}
	
	
}


\item 
(3 小题，每小题 12 分，满分 36 分)
\begin{enumerate}
	%\renewcommand{\labelenumi}{\arabic{enumi}.}
	% A(\Alph) a(\alph) I(\Roman) i(\roman) 1(\arabic)
	%设定全局标号series=example	%引用全局变量resume=example
	%[topsep=-0.3em,parsep=-0.3em,itemsep=-0.3em,partopsep=-0.3em]
	%可使用leftmargin调整列表环境左边的空白长度 [leftmargin=0em]
	\item
设 $f(x)$ 具有二阶连续导数, $f(0)=1, f^{\prime}(0)=1$, 且曲线积分 $\int_{L}\left(e^{x} \sin y+2 y f^{\prime}(x)+2 x y\right) d x+\left(f^{\prime}(x)+f(x)+2 x+e^{x} \cos y\right) d y$ 与路径无关。
\begin{enumerate}
	%\renewcommand{\labelenumi}{\arabic{enumi}.}
	% A(\Alph) a(\alph) I(\Roman) i(\roman) 1(\arabic)
	%设定全局标号series=example	%引用全局变量resume=example
	%[topsep=-0.3em,parsep=-0.3em,itemsep=-0.3em,partopsep=-0.3em]
	%可使用leftmargin调整列表环境左边的空白长度 [leftmargin=0em]
	\item
	求 $f(x)$。
	\item 
	当 $L$ 是从 $(0,0)$ 沿曲线 $y=x^{4}$ 到 $(1,1)$ 的有向曲线段时, 求该曲线积分的值。
	
	
	
	
\end{enumerate}

\item 
将函数 $y=x \arctan x-\frac{1}{2} \ln \left(1+x^{2}\right)$ 在 $x=0$ 处展成泰勒级数,
\\并求收敛域及 $ \sum\limits_{n=0}^{\infty} \frac{(-1)^{n}}{(2 n+1)(2 n+2)} $。
	
\item 
将函数 $f(x)=\left\{\begin{array}{ll}\pi & -\pi \leqslant x \leqslant 0 \\ \pi-x & 0 \leqslant x \leqslant \pi\end{array}\right.$展开成傅里叶级数 (说明收敛情况)，并求$\sum\limits_{n=1}^{\infty} \frac{1}{(2 n-1)^{2}}$。


	
\end{enumerate}


\banswer{
	\begin{enumerate}
		%\renewcommand{\labelenumi}{\arabic{enumi}.}
		% A(\Alph) a(\alph) I(\Roman) i(\roman) 1(\arabic)
		%设定全局标号series=example	%引用全局变量resume=example
		%[topsep=-0.3em,parsep=-0.3em,itemsep=-0.3em,partopsep=-0.3em]
		%可使用leftmargin调整列表环境左边的空白长度 [leftmargin=0em]
		\item
		\begin{enumerate}
			%\renewcommand{\labelenumi}{\arabic{enumi}.}
			% A(\Alph) a(\alph) I(\Roman) i(\roman) 1(\arabic)
			%设定全局标号series=example	%引用全局变量resume=example
			%[topsep=-0.3em,parsep=-0.3em,itemsep=-0.3em,partopsep=-0.3em]
			%可使用leftmargin调整列表环境左边的空白长度 [leftmargin=0em]
			\item
			$ f(x)=e^{x}-x^{2} $
			\item 
			$ e(2+\sin 1)-1 $
			
		\end{enumerate}
		\item 
		$y(x)=\sum\limits_{n=0}^{\infty} \frac{(-1)^{n} x^{2 n+2}}{(2 n+1)(2 n+2)}$，收敛域为$ |x| \leq 1 $，$ \sum\limits_{n=0}^{\infty} \frac{(-1)^{n}}{(2 n+1)(2 n+2)} =\frac{\pi}{4}- \frac{ 1 }{ 2 } \ln 2$

		\par\vspace{1mm}\noindent
		\textbf{注：}{\kaishu
		原题面印作 $y=\arctan x-\frac12\ln\left(1+x^{2}\right)$，脱落了首项的因子 $x$。
		事实上 $\left[x\arctan x-\frac12\ln\left(1+x^{2}\right)\right]^{\prime}=\arctan x$，
		逐项积分 $\arctan x=\sum\limits_{n=0}^{\infty}\frac{(-1)^{n}x^{2n+1}}{2n+1}$ 才得到上列答案；
		而 $\arctan x-\frac12\ln\left(1+x^{2}\right)$ 的展开式是
		$\sum\limits_{n=0}^{\infty}(-1)^{n}\left[\frac{x^{2n+1}}{2n+1}-\frac{x^{2n+2}}{2n+2}\right]$，
		与答案不符（如 $x=0.5$ 时前者为 $0.3521$、后者为 $0.1203$）。故已据答案更正题面。
		两式在 $x=1$ 处恰好同值，这正是所求数项级数之和 $\frac{\pi}{4}-\frac12\ln2$。
		}\par

		\item 
		$ f(x)=\frac{3}{4} \pi+\sum\limits_{n=1}^{\infty}\left[\frac{1-(-1)^{n}}{n^{2} \pi} \cos n x+\frac{(-1)^{n}}{n} \sin n x\right] \quad(-\pi<x<\pi)$，在$ x=\pm \pi $处，级数收敛于$ \frac{\pi}{2} $，$\sum\limits_{n=1}^{\infty} \frac{1}{(2 n-1)^{2}}=\frac{\pi^{2}}{8}$
		
		
		
		
	\end{enumerate}
	
	
}

\newpage
\item 
(2 小题，每小题 12 分，共 24 分)
\begin{enumerate}
	%\renewcommand{\labelenumi}{\arabic{enumi}.}
	% A(\Alph) a(\alph) I(\Roman) i(\roman) 1(\arabic)
	%设定全局标号series=example	%引用全局变量resume=example
	%[topsep=-0.3em,parsep=-0.3em,itemsep=-0.3em,partopsep=-0.3em]
	%可使用leftmargin调整列表环境左边的空白长度 [leftmargin=0em]
	\item
设 $f(x)$ 在 $[0,1]$ 上连续，在 $(0,1)$ 上可导, $f(0)=0, f(1)=2$, 证明:
\begin{enumerate}
	%\renewcommand{\labelenumi}{\arabic{enumi}.}
	% A(\Alph) a(\alph) I(\Roman) i(\roman) 1(\arabic)
	%设定全局标号series=example	%引用全局变量resume=example
	%[topsep=-0.3em,parsep=-0.3em,itemsep=-0.3em,partopsep=-0.3em]
	%可使用leftmargin调整列表环境左边的空白长度 [leftmargin=0em]
	\item
 存在 $\xi \in(0,1)$ 使 $f(\xi)=1$。
\item 
存在 $0<x_{1}<x_{2}<1$, 使 $\frac{1}{f^{\prime}\left(x_{1}\right)}+\frac{1}{f^{\prime}\left(x_{2}\right)}=1$。
\end{enumerate}
	
	
\item 
\begin{enumerate}
	%\renewcommand{\labelenumi}{\arabic{enumi}.}
	% A(\Alph) a(\alph) I(\Roman) i(\roman) 1(\arabic)
	%设定全局标号series=example	%引用全局变量resume=example
	%[topsep=-0.3em,parsep=-0.3em,itemsep=-0.3em,partopsep=-0.3em]
	%可使用leftmargin调整列表环境左边的空白长度 [leftmargin=0em]
	\item
	求 $F(x)=\int_{0}^{a}|t-x| d t$ (常数 $a>0$ ) 在 $[0, a]$ 上的最小值。
	\item 
	设 $f(x)$ 在 $[0, a](a>0)$ 连续，并且 $\int_{0}^{a} f(x) d x=0, \int_{0}^{u} x f(x) d x=1$, 求证: 存 在一点 $x_{0} \in[0, a]$ 使得 $\left|f\left(x_{0}\right)\right| \geqslant \frac{4}{a^{2}}$。
	
	
\end{enumerate}


	
	
	
	
\end{enumerate}


\banswer{
	\begin{enumerate}
		%\renewcommand{\labelenumi}{\arabic{enumi}.}
		% A(\Alph) a(\alph) I(\Roman) i(\roman) 1(\arabic)
		%设定全局标号series=example	%引用全局变量resume=example
		%[topsep=-0.3em,parsep=-0.3em,itemsep=-0.3em,partopsep=-0.3em]
		%可使用leftmargin调整列表环境左边的空白长度 [leftmargin=0em]
		\item
		\begin{enumerate}
			%\renewcommand{\labelenumi}{\arabic{enumi}.}
			% A(\Alph) a(\alph) I(\Roman) i(\roman) 1(\arabic)
			%设定全局标号series=example	%引用全局变量resume=example
			%[topsep=-0.3em,parsep=-0.3em,itemsep=-0.3em,partopsep=-0.3em]
			%可使用leftmargin调整列表环境左边的空白长度 [leftmargin=0em]
			\item
			由 $f(0)=0<1<2=f(1)$ 及 $f$ 在 $[0,1]$ 上连续，根据连续函数的介值定理，存在 $\xi \in(0,1)$ 使得 $f(\xi)=1$。
			\item 
			由(1)知存在 $\xi \in(0,1)$ 使 $f(\xi)=1$。在 $[0,\xi]$ 上对 $f$ 应用拉格朗日中值定理，存在 $x_{1} \in(0, \xi)$ 使得 $f^{\prime}(x_{1})=\frac{f(\xi)-f(0)}{\xi-0}=\frac{1}{\xi}$。在 $[\xi, 1]$ 上再应用拉格朗日中值定理，存在 $x_{2} \in(\xi, 1)$ 使得 $f^{\prime}(x_{2})=\frac{f(1)-f(\xi)}{1-\xi}=\frac{1}{1-\xi}$。于是 $\frac{1}{f^{\prime}(x_{1})}+\frac{1}{f^{\prime}(x_{2})}= \xi+(1-\xi)=1$，且 $0<x_{1}<\xi<x_{2}<1$。
			
		\end{enumerate}
		\item 
		\begin{enumerate}
			%\renewcommand{\labelenumi}{\arabic{enumi}.}
			% A(\Alph) a(\alph) I(\Roman) i(\roman) 1(\arabic)
			%设定全局标号series=example	%引用全局变量resume=example
			%[topsep=-0.3em,parsep=-0.3em,itemsep=-0.3em,partopsep=-0.3em]
			%可使用leftmargin调整列表环境左边的空白长度 [leftmargin=0em]
			\item
			$ x=\frac{a}{2} $时，$ F(x) $有最小值$ \frac{a^{2}}{4} $
			\item 
			题设中 $\int_{0}^{u} x f(x) d x=1$ 的 $u$ 未给出定义，属印刷遗漏；下在 $u \in[0, a]$ 的自然假定下完成证明。\\
			令 $M=\max_{[0, a]}|f(x)|$。若结论不真，则 $M<\frac{4}{a^{2}}$。\\
			定义 $h_{u}(x)=\begin{cases}x, & 0 \le x \le u,\\ 0, & u<x \le a.\end{cases}$ 由条件 $\int_{0}^{a} f(x) d x=0$ 及 $\int_{0}^{u} x f(x) d x=1$ 知
			\[
			\int_{0}^{a} h_{u}(x) f(x) d x=\int_{0}^{u} x f(x) d x=1.
			\]
			对任意常数 $c$ 均有 $\int_{0}^{a}(h_{u}(x)-c) f(x) d x=1-c \int_{0}^{a} f(x) d x=1$。选取适当的 $c$ 使 $\int_{0}^{a}|h_{u}(x)-c| d x$ 取得小值：
			\begin{itemize}
				\item 若 $u \le a / 2$，取 $c=0$，则 $\displaystyle\int_{0}^{a}|h_{u}(x)-c| d x=\int_{0}^{u} x d x=\frac{u^{2}}{2} \le \frac{a^{2}}{8}<\frac{a^{2}}{4}$。
				\item 若 $u \ge a / 2$，取 $c=u-a / 2\;( \ge 0)$，计算得
				\[
				\int_{0}^{a}|h_{u}(x)-c| d x
				=\int_{0}^{u}|x-c| d x+\int_{u}^{a} c d x
				=c^{2}+\frac{u^{2}}{2}-c u+c(a-u)
				=\frac{a^{2}}{4}-\frac{(a-u)^{2}}{2} \le \frac{a^{2}}{4}.
				\]
			\end{itemize}
			总之 $\int_{0}^{a}|h_{u}-c| d x \le \frac{a^{2}}{4}$。于是
			\[
			1=\left|\int_{0}^{a}(h_{u}-c) f d x\right|
			\le M \int_{0}^{a}|h_{u}-c| d x
			\le M \cdot \frac{a^{2}}{4}
			<\frac{4}{a^{2}} \cdot \frac{a^{2}}{4}=1,
			\]
			矛盾。故必存在 $x_{0} \in[0, a]$ 使得 $\left|f\left(x_{0}\right)\right| \geqslant \frac{4}{a^{2}}$。\\[2pt]
			{\small(注：原题未定义 $u$，严格说来命题不完备。)}
			
		\end{enumerate}
		
		
		
		
		
	\end{enumerate}
	
	
}


	
\end{enumerate}


